\SourceLeaf{284}%
\begin{center}\begin{tabular}{@{}r c@{}}
8.1 & $-$\\
9.2 & $+$\\
12.5 & $-$\\
14.7 & $+$\\
16.11 & $-$\\
17.4 & $-$\\
18.3 & $+$\\
19.10 & $+$\\
20.15 & $+$\\
21.6 & $-$\\
22.13 & $-$\\
\end{tabular}\end{center}

Es ist klar, dass, indem man von einem Punkte vor (1) ausgeht und in der
Curve die Punkte 1, 2, 3 etc. durchläuft, sich zum erstenmale bei dem
Schnitt 8 eine Figur bildet, die wir $a$ nennen wollen, diese geht von 1
bis 8; offenbar liegen der Theil der ganzen Linie zunächst vor 1 und
[der Theil] nach 8 auf der $-$Seite von $a$.

Indem man so weiter fortgeht, bildet sich eine neue zusammenhängende
Curve, die vom Anfangspunkt bis 1 (dann den vorigen Weg 2, 3 etc.
verlassend) auf 9, 10, 11 etc. fortgeht; bei 9 gelangt man auf die
$+$Seite von $a$, bei 12 auf die $-$, bei 14 auf die $+$Seite von $a$;
bei 16 hingegen trifft die Curve sich selbst und schliesst eine zweite
Figur, die von 11 durch 12, 13, 14, 15 bis 16 geht, diese heisse $b$;
offenbar liegt der Theil der Curve zunächst vor 1, dann das ganze
Stück von 8 bis 11, endlich der zunächst an 16 grenzende Theil auf
der $-$Seite von $b$. Nachdem jetzt auch $b$ abgesondert ist, trifft die
Curve in 17 wieder $a$ und kommt auf die $-$Seite, in 18 abermals, indem
sie auf die $+$Seite gelangt; in 19 schliesst sich die $3^{\mathrm{te}}$
Figur (da 10 noch nicht in einem ausgeschlossenen Theil liegt), welche
$c$ heissen soll und von 10 bis 11 $=$ 16, dann von 16 bis 19 geht, und
wenn diese wieder abgesondert wird, liegt das übrig bleibende Stück
so wie das nächst folgende auf der $+$Seite von $c$; bei 20 gelangt
man auf die $+$Seite von $b$, bei 21 auf die $-$Seite von $a$, bei 22
auf die $-$Seite von $b$ und, indem man auf den Anfangspunkt (vor 1)
zurückkommt, schliesst sich die $4^{\mathrm{te}}$ Figur, die $d$
heisst.
\par\Needspace{30\baselineskip}
Die Lage gegen die schon beschriebenen Curven ergibt sich also
